\section{Why the transpose appears in the field equation}
The transpose is not decorative notation. It makes the dimensions of an
operation visible. If there are $n$ cells and $m$ available action directions,
$S$ has $n$ rows and $m$ columns. A quantity vector $q$ has $m$ entries,
so $Sq$ has $n$ entries: one stock change for each cell. The marginal
$\mu$ has $n$ entries. Its transpose is a row, and $\mu^TSq$ is a
single number. That number is the tangent potential change produced by
the complete quantity vector.

Associativity lets us calculate the same scalar in two stages. We can first
form the physical change $Sq$ and then value its direction with $\mu^T$.
Or we can first form the row $\mu^TS$, one marginal change per action
direction, and then multiply by the quantities. Since
$\mu^TS=(S^T\mu)^T$, defining $f=-S^T\mu$ gives
\begin{equation}
 -\mu^TSq=f^Tq.
\end{equation}
This is the whole sign convention in one line. A positive entry of $f$
means that a small positive quantity in that declared direction reduces
the potential. It does not authorize an action or select its quantity.

For the chain matrix earlier in the chapter, let
$\mu=(2,0,-2)^T$. Then $S^T\mu=(-2,-2)^T$, so both edge forces
are two. For $q=(3,2)^T$, the tangent benefit is ten. Directly,
$Sq=(-3,1,2)^T$, and $-\mu^TSq=-(-6-4)=10$.
This checks the matrix arithmetic, not the exact finite value. The latter
must still subtract the quadratic curvature contribution.

\section{Addition of vectors does not establish simultaneous permission}
Suppose a proposed group has three children. Writing their increments as
columns makes their sum easy to inspect. A coordinate that is positive in
one child and negative in another may cancel in the aggregate. Algebra
records that cancellation exactly. Physical interpretation requires more.
Did the middle tank hold the resource long enough to make the outgoing
delivery? Was an intermediate conversion required? Could the two devices
operate concurrently? These questions are not answered by a zero in the
aggregate vector.

The distinction is similar to adding deposits and withdrawals in a stock
record. A zero net daily change does not prove that the stock was never
empty during the day. A constant-rate simultaneous model can make a
specific path declaration; an endpoint-only model leaves intermediate
behavior unspecified. A chronological model needs actual ordered states.
All three can share endpoints while representing different physical claims.

\section{Weighted length and the quadratic form}
The expression $v^THv$ looks more complicated than it is. With diagonal
$H$, first multiply each entry $v_i$ by $1/\sigma_i^2$, then multiply
again by $v_i$ and sum:
\begin{equation}
 v^THv=\sum_i\frac{v_i^2}{\sigma_i^2}.
\end{equation}
It is a squared length measured in the declared scales. Its value cannot
be negative, and it vanishes only when every entry vanishes. Reversing a
vector leaves this squared length unchanged, although it reverses the
linear directional term $\mu^Tv$.

This difference between the linear and quadratic terms explains a common
finite-action surprise. A small movement toward the reference can be
helpful, but extending the same movement far enough overshoots. The linear
part grows with quantity; the curvature term grows with its square. A
longer movement is therefore not just several copies of the original
marginal improvement at an unchanged state.

For the chain example with equal scales two and increment $(-3,1,2)$,
the squared weighted length is $(9+1+4)/4=7/2$. Its half is $7/4$.
If the marginals above came from a compatible Gaussian declaration,
the exact finite field benefit would be $10-7/4=33/4$. For instance,
state $(18,10,2)$ and reference $(10,10,10)$ provide precisely those
marginals. The endpoint is $(15,11,4)$, with potential $31/4$ rather
than the starting sixteen. Their difference is again $33/4$.

\section{Coordinate changes and model changes}
Reordering the list of cells changes the written vector, but not the
physical account if every associated object is reordered consistently.
The state, reference, scales, incidence rows and ownership identifiers
must all follow the same permutation. Reordering just the state can turn
a formally well-shaped matrix multiplication into the wrong physical
calculation without producing a syntax error.

Changing litres to millilitres is another consistent coordinate change,
already checked above. Replacing a scale because one cell is judged more
important is not. Nor is combining two distinct carriers into one scalar
because they share a convenient numeric range. The first changes units;
the second changes valuation; the third changes the represented world.
Software needs identities that make these differences visible.

\section{Practice: follow one symbol through every layer}
Choose one edge quantity $q_e$. Identify its stock unit, its incidence
column, the cells it changes, the constraints it consumes, and the
corresponding force entry. Then identify the finite curvature term in
which it appears alone and any cross terms in which it meets another
quantity. Finally ask whether a named owner exists. The matrix will not
supply that final answer.

This exercise is worth doing before proceeding to proofs. Many later
errors are not failures of advanced calculus. They are failures to carry
one symbol's meaning consistently through a chain of otherwise correct
operations. With this dimensional foundation, the next chapter can
introduce the potential without asking the reader to trust opaque notation.
